\documentclass[11pt,a4paper]{article}
\usepackage{../shared/hanzocover}
\usepackage[utf8]{inputenc}
\usepackage[T1]{fontenc}
% Shared preamble for the compute-market series.
% One style, one vocabulary, inputted by every paper so the series reads as one voice.
\usepackage[margin=1in]{geometry}
\usepackage{booktabs}
\usepackage{amsmath}
\usepackage{amssymb}
\usepackage{amsthm}
\usepackage{graphicx}
\usepackage{xcolor}
\usepackage{hyperref}
\hypersetup{colorlinks=true, linkcolor=blue, urlcolor=blue, citecolor=blue}
\usepackage{enumitem}
\usepackage{listings}
\usepackage{array}
\usepackage{tabularx}
\usepackage{siunitx}
\sisetup{group-separator={,}, group-minimum-digits=4}

\definecolor{kw}{rgb}{0.13,0.29,0.53}
\definecolor{cmt}{rgb}{0.40,0.40,0.40}
\definecolor{str}{rgb}{0.16,0.50,0.16}
\lstset{
  basicstyle=\ttfamily\footnotesize,
  keywordstyle=\color{kw}\bfseries,
  commentstyle=\color{cmt}\itshape,
  stringstyle=\color{str},
  showstringspaces=false,
  breaklines=true,
  columns=fullflexible,
  frame=single,
  framesep=4pt,
}

\newtheorem{definition}{Definition}
\newtheorem{proposition}{Proposition}
\newtheorem{remark}{Remark}

\newcommand{\code}[1]{\texttt{#1}}
\newcommand{\repo}[1]{\texttt{#1}}

% Series vocabulary — used identically in every paper.
\newcommand{\job}{\ensuremath{J}}
\newcommand{\bundle}{\ensuremath{R}}
\newcommand{\cost}{\ensuremath{\mathcal{C}}}

\tolerance=800
\emergencystretch=3em


\title{Open Weights, Bought Data, Rented Gradients:\\Training as a Market Workload}

\author{
Hanzo AI Research\thanks{Corresponding author: research@hanzo.ai} \\
\textit{Hanzo Industries Inc (Techstars '17)} \\
Los Angeles, California \\
\texttt{https://hanzo.ai}
}

\date{2026}

\begin{document}
\hanzocoverpage

\begin{abstract}
A compute market that sells only inference is selling the cheap half. Training
is where the capital goes, where the differentiation lives, and where a
decentralized provider network has a structural advantage that inference does not
give it --- training tolerates latency, tolerates heterogeneity, and is
checkpointable, which means it can use exactly the supply that interactive serving
cannot.

It also has a verification problem that is worse than inference's and a data
problem that inference does not have at all. This paper treats all three. We
describe why open-weight models are the only viable substrate for a decentralized
training market, what the training workload actually looks like when decomposed
into market-tradable units, how gradient work can be checked without recomputing
it, and how data becomes a purchasable input with provenance that survives into
the resulting weights.

The enabling development is that confidential computing on current accelerators
has become close to free for the workloads that matter. That changes the design:
a provider can be given proprietary weights and licensed data, prove it ran the
declared computation inside an attested boundary, and be unable to exfiltrate
either. Training on rented hardware stops requiring that the renter be trusted.
\end{abstract}

\tableofcontents
\newpage

\section{Why open weights are the substrate}

A decentralized training market cannot be built on closed models, and the reason
is structural rather than ideological. A closed model reachable only through a
vendor's API cannot be fine-tuned on a third party's accelerator, because the
weights never leave the vendor. Whatever a market does with such a model, it is
reselling an API, and the provider network is decoration.

Open-weight models invert this. The weights are an artifact anyone can hold, so
any provider with sufficient memory is a candidate supplier, and the market has
something real to allocate. Everything in this paper assumes open weights for
that reason.

This produces a natural division of labour with a routing layer over closed
frontier models. The closed family is where the best answer often comes from
today, and routing across it is a real product. The open family is where a
decentralized network can actually do work --- serve on arbitrary hardware, be
fine-tuned per tenant, be distilled, be quantized, be improved by anyone. A market
that carries both routes intelligence at the top and allocates genuine compute
underneath.

\section{Decomposing the workload}

Training is not one job. It decomposes into units with sharply different market
properties, and a market that treats them alike will misprice all of them.

\begin{table}[h]
\centering
\begin{tabularx}{\textwidth}{lXlll}
\toprule
Unit & What it is & Latency & Parallel & Verifiable \\
\midrule
Adapter fit & low-rank delta on frozen base & tolerant & across tenants & partly \\
Preference optimisation & direct preference and its relatives & tolerant & across pairs & partly \\
Full fine-tune & all weights, sharded & tolerant & across shards & hard \\
Distillation & student fits teacher outputs & tolerant & across batches & yes \\
Synthetic generation & teacher produces training data & tolerant & embarrassingly & yes \\
Evaluation & scored inference on a fixed set & tolerant & embarrassingly & yes \\
Quantization & post-training weight compression & tolerant & across tensors & yes \\
\bottomrule

\end{tabularx}
\caption{Training units and their market properties.}
\end{table}

Three of these are worth singling out because they are the easy money and are
routinely overlooked.

\paragraph{Synthetic generation is inference.} Producing training data from a
teacher model is a batch inference workload with no latency constraint and no
interactivity. It is the single best fit for cheap heterogeneous supply in the
entire stack, and it verifies under exactly the determinism discipline the
companion paper describes --- pin the seed, pin the numerics, and the output is
reproducible.

\paragraph{Evaluation is inference on a fixed set.} Scoring a candidate model
against a benchmark is deterministic given a pinned harness. A market can sell
evaluation as a service with byte-equality verification and no trust
requirements, which also gives it something valuable: an independent measurement
of every model in its own catalogue.

\paragraph{Adapter fitting parallelises across tenants, not within.} Ten thousand
tenants each fitting a small adapter is ten thousand independent jobs. That is a
better fit for a heterogeneous network than one large job needing tight
interconnect, and it is the workload personalization actually generates.

By contrast, a full fine-tune of a large model needs high-bandwidth interconnect
between shards, which is precisely what a geographically dispersed network does
not have. Honest market design puts full pretraining and large full fine-tunes in
the datacentre and takes the rest.

\section{Checking gradient work}

Training's verification problem is harder than inference's. An inference result
can at least be spot-checked algebraically against the weights that produced it.
A gradient step's output is a weight update, and a provider that returns a
plausible but wrong update degrades the model subtly over many steps rather than
failing visibly once.

Four mechanisms are available, and they stack.

\paragraph{Loss trajectory as a cheap signal.} The buyer holds out a small
validation set and evaluates the returned checkpoint. A provider submitting
garbage is caught immediately; a provider submitting a slightly-wrong update is
not. This is a liveness check, not a correctness check, and should be described as
such.

\paragraph{Recomputation of sampled steps.} The buyer recomputes a random fraction
of steps from the checkpoint and the batch. This requires the step to be
deterministic --- pinned numerics, pinned data order, pinned dropout seed ---
which is the same discipline as everywhere else in this series. Cost is the
sampling rate; soundness is the sampling rate. With stake at risk, a low rate is
adequate, and it is the workhorse mechanism.

\paragraph{Algebraic checks on the dominant operation.} The backward pass is
dominated by matrix products, and those are checkable in quadratic rather than
cubic time by randomised verification, at well under a percent of the cost of the
product for the dimensions in question. This checks the expensive part and is
silent about the rest.

\paragraph{Attested execution.} Discussed next, and increasingly the answer.

\section{Confidential computing changes the design}

The development that reorganises this problem is that hardware confidential
computing on current accelerators has become cheap enough to use by default.

For compute-bound large-model work --- which is what training and large-batch
inference are --- the overhead of running inside an encrypted memory boundary is
small, because the cost is dominated by arithmetic rather than by transfers
across the boundary. Current-generation hardware with encrypted device links
narrows the remaining gap further by removing the staging copies that earlier
designs required. Small, transfer-bound workloads still pay meaningfully; the
workloads a compute market sells largely do not.

We classify providers by what their hardware can actually attest, in four tiers,
with all attestation performed locally rather than against a vendor cloud service
--- a market that cannot verify without an external dependency has not removed the
trusted party, it has renamed it.

\begin{table}[h]
\centering
\begin{tabular}{llll}
\toprule
Tier & Hardware & What it attests & Trust band \\
\midrule
1 & accelerator-native confidential mode & device, firmware, encrypted memory & 90--100 \\
2 & confidential virtual machine plus accelerator & host measurement plus device & 70--89 \\
3 & device secure enclave with an inference engine & enclave measurement & 50--69 \\
4 & no hardware support & nothing; stake substitutes & 10--49 \\
\bottomrule
\end{tabular}
\caption{Provider tiers by attestable capability. An attestation carries the
device identity, model, whether confidential mode and encrypted device links are
enabled, the partitioning configuration, and driver and firmware versions.}
\end{table}

\subsection{What this makes possible}

\paragraph{Training on weights the provider may not keep.} A model owner ships
proprietary weights into an attested boundary. The provider supplies
accelerators and receives payment; it does not receive the weights in any form it
can extract. This is the difference between a market that can only train public
models and one that can train anything.

\paragraph{Training on licensed data.} The same argument applies to the data,
and it matters more, because data licences are typically far more restrictive
than model licences. A corpus that may not be copied can still be trained on
inside a boundary that can prove it did not leave.

\paragraph{Verification for nondeterministic work.} This is the deepest
consequence. Byte-equality verification covers bit-reproducible work and stops
there. Attested execution covers the rest: it proves the declared computation ran
on declared hardware with declared inputs, which is exactly the class of claim
that no algebraic check can make. At a small percentage overhead this is
dramatically cheaper than replication, and it applies where the cheap checks
cannot reach.

\subsection{What it does not do}

Attestation moves the trust root to a silicon vendor's certificate chain. That is
a real and different assumption, not an elimination of assumptions, and the
history of side-channel and fault attacks against such boundaries argues against
using attestation as the sole mechanism for the highest-value work. Used together
with sampled recomputation and algebraic checks it is excellent; used alone at
tier 1 it is a strong claim resting on one vendor's supply chain.

\section{Data as a purchasable input}

Inference buys compute. Training buys compute and data, and the data side is the
one most markets have no mechanism for at all.

\paragraph{The paradox.} Data's value cannot be assessed without inspecting it,
and inspecting it is most of what a buyer wanted to pay for. The mechanisms that
resolve this are all forms of buying a \emph{function} of the data rather than the
data itself.

\paragraph{Train-without-transfer.} The buyer sends the model into the attested
boundary; the seller supplies data that never leaves it; the buyer receives the
updated weights. The seller is paid for the improvement rather than for the
corpus, and retains it. This is the mechanism confidential computing unlocks, and
it is the reason the two halves of this paper belong together.

\paragraph{Evaluation before purchase.} The seller permits a scored evaluation
run against a small sample, producing a measured improvement figure that the buyer
can use to price the whole. This requires trusting that the sample represents the
corpus, which is a bounded and familiar commercial risk.

\paragraph{Provenance that survives into the weights.} A training run should
record which corpora contributed, under what licence, in what proportion, with
content-addressed identifiers --- carried with the resulting checkpoint. This is
becoming a compliance requirement rather than a nicety, and a market that records
it structurally is worth more to any buyer who will ever be asked where their
model came from. It also creates the ledger against which contributors can be paid
when a model is later commercialised, which is the only credible basis for
recurring data compensation.

\paragraph{Contribution measurement.} Attributing model quality to individual data
sources is genuinely unsolved at scale --- influence functions and leave-one-out
retraining are both far too expensive for large corpora. We do not claim
otherwise. Proportional payment by tokens contributed, weighted by a measured
improvement on held-out evaluation for the corpus as a whole, is crude, cheap and
implementable, and it is better than the alternative of paying nothing.

\section{What exists}

The training side of our stack is open source and in production use.

A fine-tuning framework carrying low-rank and quantized low-rank adaptation,
direct preference optimisation and its group-relative, odds-ratio and
Kahneman--Tversky variants, proximal policy optimisation, sharded data-parallel
and pipeline execution, sample packing and fused attention --- the standard corpus
of open-weight training methods, maintained as a non-profit research project.

A training crate inside the inference engine implementing trainable low-rank
layers over frozen bases, an optimiser, distillation and adapter serialisation
whose tensor layout matches the inference-side merge, exposed over the serving API
as a session-oriented loop: create a client against a base model, accumulate
gradients, step, sample from the current weights, save the adapter.

Provider classification and local attestation verification for the four tiers
above, including accelerator attestation parsing.

And the pricing and routing layer that quotes all of it, described in the
companion paper on the market mechanism.

\section{Sequencing}

Synthetic generation and evaluation first: both are inference, both verify under
byte equality, both tolerate arbitrary latency, and both have immediate demand
from anyone training anything. Adapter fitting next, because personalization
generates it in volume and it parallelises across tenants rather than within a
job. Distillation and quantization after that. Attested training on proprietary
weights and licensed data once tier-1 supply is a meaningful share of the network
--- that is the highest-value workload here, and it is the one that requires the
hardware story to be real rather than aspirational.

Large full fine-tunes and pretraining stay in the datacentre. Saying so plainly is
worth more than a claim that will be tested and fail.

\end{document}
